% !TeX root = ../main.tex
基于Wi-Fi CSI的室内动作识别是当前无线感知与智能感知领域的热点研究课题，论文针对Wi-Fi CSI动作识别中类别持续扩展、旧类数据难以完整保留以及新类样本不足条件下的模型持续更新问题开展研究工作，选题具备一定的理论价值。

论文的主要工作包括以下三点：1、针对新增类样本相对充足的类增量学习场景，提出回放——扩展——巩固框架；2、针对新增类样本极少的少样本增量学习场景，提出基于时序增强与图注意力的学习方案；3、在XRF和WiAR数据集上进行了对比实验和消融实验，验证了所提方法的有效性。

论文写作认真，结构清晰，论述完整。论文工作表明作者基本掌握软件工程学科的基础理论和专业知识，具有一定的独立从事科研工作的能力。

答辩中讲述清晰，回答问题基本正确，经答辩委员会讨论和无记名投票表决，一致同意通过学位论文答辩，并一致建议授予罗弘林软件工程专业硕士学位。
